% Independent proof insert; source for slope is Bruguières--Virelizier,
% On the center of fusion categories, Pacific J. Math. 264 (2013), 1--30, Section 4C.
% Source for the ring parameterization is Larson, arXiv:1401.1879v2, Prop. 2.1.
\subsection{Removing the multiplicity bound in the non-self-dual pivotal case}

\begin{lemma}\label{lem:rank4-nonselfdual-spherical}
Every pivotal structure on a nonpointed rank-four fusion category with a
non-self-dual simple object is spherical.
\end{lemma}
\begin{proof}
Write the simple objects as $1,X,Y,Z$, with $X^*=Z$ and $Y^*=Y$.
The slope operator of a pivotal fusion category is the monoidal natural
endomorphism of the identity whose scalar on a simple object $V$ is
$\operatorname{sl}(V)=\dim_l(V)/\dim_r(V)$; see
\cite[Section~4C]{BVCenter}.
Consequently $N_{UV}^W>0$ implies
$\operatorname{sl}(W)=\operatorname{sl}(U)\operatorname{sl}(V)$.
The slopes of $1,Y$ are one, and those of $X,Z$ are $s,s^{-1}$.
We prove $s=1$.

Use the general rank-four non-self-dual multiplication table
\[
\begin{aligned}
 X^2&=pX+lY+cZ,& X Y&=qX+kY+lZ,\\
 Y^2&=1+kX+eY+kZ,& XZ&=1+pX+qY+pZ,\\
 Z^2&=cX+lY+pZ,& YZ&=lX+kY+qZ,
\end{aligned}
\]
where $c,e,k,l,p,q$ are nonnegative integers. Associativity gives
\[
 kl+lc=lp+kq,\qquad kp+le+kc=2lq+k^2,
\]
\[
 l^2+c^2=1+q^2+p^2,\qquad l^2+k^2+q^2=1+2pk+qe.
\]
If $s\ne1$, the occurrences of $X$ in $X^2$ and of $Y$ in $XY$ force
$p=k=0$. The first associativity equation then gives $lc=0$. If $l=0$,
the third equation gives $(c-q)(c+q)=1$, hence $c=1,q=0$, contrary to
the fourth equation. Therefore $l>0$ and $c=0$. The third equation now
gives $(l-q)(l+q)=1$, hence $l=1,q=0$; the second then gives $e=0$.
The resulting table is the pointed ring of $\mathbb Z/4\mathbb Z$,
contrary to the hypothesis. Thus every slope is one, which is precisely
sphericity.
\end{proof}

\begin{proposition}\label{prop:nonselfdual-pivotal-pu}
A rank-four non-self-dual fusion ring admits a pivotal categorification
if and only if it admits a pseudo-unitary categorification.
\end{proposition}
\begin{proof}
Only the forward implication requires proof. If the ring is pointed,
it admits a unitary model. Otherwise the preceding lemma makes any
pivotal realization spherical.

Let $R$ be the ring. Its four complex characters are permuted by the
involution $\chi\mapsto(b\mapsto\chi(b^*))$. On $R\otimes_{\mathbb Z}\mathbb C$
the involution has trace two, since it fixes the two basis elements
$1,Y$ and exchanges $X,Z$. On the primitive idempotent basis its trace
is the number of fixed characters. Hence precisely two characters are
fixed by duality. These are exactly the real-valued characters, because
characters of a commutative fusion ring satisfy
$\chi(b^*)=\overline{\chi(b)}$.

One of these real characters is the Frobenius--Perron character $d$.
If all its values are rational, they are integers, and every fusion
categorification is weakly integral and therefore pseudo-unitary.
Otherwise the Galois orbit of $d$ has at least two members. It consists
entirely of characters fixed by duality, since $d(b)=d(b^*)$ is preserved
by field automorphisms. Thus it consists of both real-valued characters.
The spherical dimension character is real-valued, so is Galois conjugate
to $d$. Conjugating the category and its spherical structure gives
positive Frobenius--Perron dimensions, hence a pseudo-unitary realization.
\end{proof}

\begin{corollary}\label{cor:nonselfdual-unbounded-classification}
With no multiplicity restriction, the rank-four non-self-dual fusion
rings admitting pivotal categorifications are exactly the six rings
\[
 \mathbb Z[\mathbb Z/4\mathbb Z],\quad
 \operatorname{NG}(\mathbb Z/3\mathbb Z,m)\ (m=0,2,3,6),\quad S(1).
\]
All six admit unitary realizations.
\end{corollary}
\begin{proof}
The proposition reduces the classification to the pseudo-unitary case.
Larson's classification gives these six rings and the additional candidate
$S(2)$. The uniform exclusion of $S(m)$ for $m\geq2$ removes that candidate.
The known unitary models are the ones already recorded in the realization
table: the cyclic group, Tambara--Yamagami, $\operatorname{Rep}(A_4)$,
Izumi--Xu, the $m=6$ near-group model, and the non-self-dual quadratic model.
\end{proof}

\begin{corollary}\label{cor:rank4-pivotal-spherical-unbounded}
Every pivotal structure on a nonpointed rank-four fusion category is
spherical. Consequently, pivotal and spherical categorifiability are
equivalent for all rank-four fusion rings, without a multiplicity bound.
\end{corollary}
\begin{proof}
In the self-dual case, left and right dimensions agree on each simple
object. The non-self-dual case is Lemma~\ref{lem:rank4-nonselfdual-spherical}.
Every pointed rank-four ring has a unitary realization, which handles
the remaining existence statement.
\end{proof}
